\begin{explanationbox}
	\textbf{\textcolor{CExplanation}{一句话直觉}}

	从左往右看，两段各自向上走，跨过分界点也不能往下掉，整条图象才会递增；递减时全部反过来。

	\medskip
	\textbf{\textcolor{CExplanation}{要检查的两件事}}
	\begin{enumerate}[leftmargin=*, itemsep=0.35em]
		\item \textbf{段内方向：} $g$ 与 $h$ 分别在自己的定义区间内严格递增（或严格递减）。
		\item \textbf{跨段衔接：} 递增时检查 $g(x_0)\le L$；递减时检查 $g(x_0)\ge L$，其中 $L=\lim_{x\to x_0^+}h(x)$。
	\end{enumerate}

	\textbf{\textcolor{CExplanation}{为什么接口可以取等？}}

	设 $x_1\le x_0<x_2$。严格递增时，$g(x_1)\le g(x_0)\le L<h(x_2)$，所以 $f(x_1)<f(x_2)$。
	右段从 $x_0$ 的右边才开始，即使左端函数值恰好等于右极限，任意真正属于右段的点仍严格高于这个极限。递减时有
	$g(x_1)\ge g(x_0)\ge L>h(x_2)$。

	\medskip
	\textbf{\textcolor{CExplanation}{图象提醒}}

	接口可以向单调方向\textbf{跳跃}，不要求连续，更不要求光滑。递增时只禁止向下跳；递减时只禁止向上跳。

	\textbf{\textcolor{CExplanation}{怎样处理 $h(x_0)$？}}

	$x_0$ 不属于右段。应先求 $L$；如果右段表达式在 $x_0$ 可代入，且它在该点可连续延拓，那么代入得到的数值恰好等于 $L$。代入只是\textbf{计算极限的方法}，并不改变 $f(x_0)=g(x_0)$。

	\textbf{\textcolor{CExplanation}{常见考法}}

	含参题先分别求出两段严格单调的参数条件，再列接口不等式，最后取交集。
\end{explanationbox}